\documentclass[12pt]{article}

% ====== Page Layout and Basic Formatting ======
\usepackage[left=1.0in,right=1.0in,top=1.0in,bottom=1.0in]{geometry}
\usepackage{setspace}
\doublespacing
\usepackage{fancyhdr}
\pagestyle{fancy}
\fancyhf{}
\fancyfoot[C]{\thepage}
\renewcommand{\headrulewidth}{0pt}

% ====== Typography and Fonts ======
\usepackage{lmodern}
\usepackage{microtype}
\usepackage[normalem]{ulem}
\usepackage[T1]{fontenc}
\usepackage[spanish,es-noquoting,es-tabla]{babel}

% ====== Section Styling ======
\usepackage{titlesec}
\usepackage[title]{appendix}
\usepackage{titling}
\pretitle{\begin{center}\large\bfseries}
\posttitle{\end{center}}
\preauthor{\begin{center}\normalsize}
\postauthor{\end{center}}
\predate{\begin{center}\normalsize}
\postdate{\end{center}}

% ====== Math Packages (mínimo para artículo jurídico) ======
\usepackage{amssymb, amsmath}

% ====== Table Packages ======
\usepackage{array, booktabs, makecell}
\usepackage[flushleft]{threeparttable}
\usepackage{tabularx}
\usepackage{tabularray}
\UseTblrLibrary{booktabs}

% ====== Figure and Caption Packages ======
\usepackage{graphicx}
\usepackage{caption}
\captionsetup{font=small, labelfont=bf, justification=justified}

% ====== List Formatting ======
\usepackage{enumitem}

% ====== Bibliography and Citation (biblatex + biber) ======
\usepackage{xurl}
\usepackage{xcolor}
\definecolor{citationcolor}{RGB}{0, 127, 255}

\usepackage[backend=biber,
            style=authoryear,
            maxcitenames=3,
            mincitenames=1,
            maxbibnames=99,
            giveninits=true,
            uniquename=false,
            uniquelist=true,
            dashed=false,
            urldate=long,
            url=true,
            natbib=true]{biblatex}
\addbibresource{../Bibliography_base.bib}

\renewcommand*{\nameyeardelim}{\addcomma\space}
\DeclareCiteCommand{\cite}
  {\usebibmacro{prenote}}
  {\usebibmacro{citeindex}%
   \printtext[bibhyperref]{\color{citationcolor}\usebibmacro{cite}}}
  {\multicitedelim}
  {\usebibmacro{postnote}}
\DeclareCiteCommand{\parencite}[\mkbibparens]
  {\usebibmacro{prenote}}
  {\usebibmacro{citeindex}%
   \printtext[bibhyperref]{\color{citationcolor}\usebibmacro{cite}}}
  {\multicitedelim}
  {\usebibmacro{postnote}}

% ====== Footnote Settings ======
\interfootnotelinepenalty=10000
\setlength{\footnotesep}{0.5cm}

% ====== URL Bleeding Fixes ======
\setcounter{biburllcpenalty}{7000}
\setcounter{biburlucpenalty}{8000}
\setcounter{biburlnumpenalty}{9000}

% ====== Hyperref (loaded second-to-last) ======
\usepackage[hidelinks, breaklinks, colorlinks=true,
            linkcolor=citationcolor, citecolor=citationcolor,
            urlcolor=citationcolor]{hyperref}

% ====== Cleveref (loaded after hyperref) ======
\usepackage[nameinlink]{cleveref}

% ====== Comandos para derecho ======
% Abreviatura de casación
\newcommand{\cas}[1]{\textit{CAS.~#1}}
% Artículo del Código Civil
\newcommand{\ccart}[1]{art.~#1~CC}

% ==========================================================

\title{%
  \textit{Lex artis}, consentimiento informado y daño médico:\\
  criterios jurisprudenciales de responsabilidad civil en el Perú%
  \thanks{El autor agradece a [mencionar si aplica].
  Las opiniones expresadas son de exclusiva responsabilidad del autor.}%
}

\author{%
  [Nombre del autor]%
  \thanks{Docente de la Facultad de Derecho, Universidad César Vallejo.
  Correo electrónico: \href{mailto:correo@ucv.edu.pe}{correo@ucv.edu.pe}.}%
}

\date{Mayo 2026}

% ==========================================================

\begin{document}

\maketitle
\thispagestyle{empty}

% ====== Resumen ======
\begin{abstract}
\noindent \singlespacing
El presente artículo analiza los criterios jurisprudenciales que emplea la judicatura peruana
para determinar la responsabilidad civil en el ámbito médico, con especial atención a tres ejes:
la \textit{lex artis} médica, el consentimiento informado y el daño médico resarcible.
A partir de un análisis dogmático del Código Civil de 1984 y de la normativa sanitaria especial
(Ley 26842 y Ley 29414), y mediante el estudio de sentencias de la Corte Suprema de Justicia
y Cortes Superiores emitidas entre 2010 y 2024, se identifican patrones y divergencias en la
aplicación de estos criterios. El artículo concluye con propuestas orientadas a fortalecer la
coherencia y predictibilidad del sistema de responsabilidad civil médica peruano.
\end{abstract}

\vspace{0.8em}
\noindent \textbf{Palabras clave:} responsabilidad civil médica, \textit{lex artis}, consentimiento
informado, daño médico, jurisprudencia peruana, Código Civil del Perú.

\vspace{0.4em}
\noindent \textbf{Área temática:} Derecho Civil — Derecho de Daños.

\vspace{1.2em}

% ====== Abstract (English) ======
\begin{otherlanguage}{english}
\begin{abstract}
\noindent \singlespacing
This article examines the jurisprudential criteria applied by Peruvian courts to determine
civil liability in the medical field, focusing on three axes: medical \textit{lex artis},
informed consent, and compensable medical harm. Through a doctrinal analysis of the 1984
Civil Code and special health legislation (Laws 26842 and 29414), and by studying Supreme
Court and Superior Court decisions from 2010 to 2024, the article identifies patterns and
divergences in the application of these standards. The article concludes with proposals aimed
at enhancing the coherence and predictability of Peru's medical civil liability system.
\end{abstract}

\noindent \textbf{Keywords:} medical civil liability, \textit{lex artis}, informed consent,
medical harm, Peruvian case law, Peruvian Civil Code.
\end{otherlanguage}

\newpage
\setcounter{page}{1}

% ==========================================================
% CUERPO DEL ARTÍCULO
% ==========================================================

% SECCIÓN 1: INTRODUCCIÓN
%
% Objetivos de esta sección:
%   1. Presentar el problema jurídico central: la falta de criterios uniformes
%      en la jurisprudencia peruana para determinar responsabilidad civil médica.
%   2. Delimitar el objeto de estudio: lex artis, consentimiento informado y daño médico
%      como los tres ejes de la responsabilidad civil en el ámbito sanitario.
%   3. Justificar la relevancia práctica y dogmática del estudio.
%   4. Anunciar la metodología: análisis dogmático + análisis jurisprudencial
%      de sentencias de la Corte Suprema y Cortes Superiores del Perú.
%   5. Presentar la estructura del artículo.
%
% Extensión estimada: 600–900 palabras

\section{Introducción}
\label{sec:introduccion}



% SECCIÓN 2: LA RESPONSABILIDAD CIVIL MÉDICA EN EL ORDENAMIENTO JURÍDICO PERUANO
%
% Objetivos de esta sección:
%   1. Enmarcar la responsabilidad civil médica en el sistema general del CC de 1984:
%      - Régimen contractual (art. 1314 CC): obligación de medios vs. de resultado.
%      - Régimen extracontractual (art. 1969 CC): culpa y antijuridicidad.
%   2. Explicar la dualidad del régimen peruano y sus consecuencias prácticas
%      (plazo de prescripción, carga de la prueba, quantum indemnizatorio).
%   3. Identificar los cuatro elementos de la responsabilidad civil:
%      antijuridicidad, factor de atribución (culpa / riesgo), nexo causal, daño.
%   4. Situar la Ley 26842 (Ley General de Salud) y la Ley 29414
%      como normativa especial que complementa el CC.
%
% Referencias clave: Espinoza Espinoza, Taboada Córdova, De Trazegnies, art. 1969-1985 CC
% Extensión estimada: 1.200–1.800 palabras

\section{La responsabilidad civil médica en el ordenamiento jurídico peruano}
\label{sec:rc_medica}



% SECCIÓN 3: LA LEX ARTIS MÉDICA — NATURALEZA, FUNCIÓN Y CRITERIOS DE DETERMINACIÓN
%
% Objetivos de esta sección:
%   1. Definir la lex artis médica como el estándar de conducta exigible al profesional
%      de salud en el ejercicio de su actividad.
%   2. Distinguir:
%      - Lex artis general: estándar de la especialidad médica (objetivado).
%      - Lex artis ad hoc: estándar ajustado a las circunstancias del caso concreto
%        (paciente, recursos disponibles, urgencia).
%   3. Analizar cómo el incumplimiento de la lex artis configura la antijuridicidad
%      y el factor de atribución (culpa profesional).
%   4. Revisar los criterios usados en la jurisprudencia comparada (España, Argentina)
%      para la determinación de la lex artis.
%   5. Identificar problemas probatorios: prueba pericial, historia clínica, protocolos.
%
% Referencias clave: Fernández Cruz, Galán Cortés, Lorenzetti, CAS. relevantes
% Extensión estimada: 1.500–2.000 palabras

\section{La \textit{lex artis} médica: naturaleza, función y criterios de determinación}
\label{sec:lex_artis}



% SECCIÓN 4: EL CONSENTIMIENTO INFORMADO — CONFIGURACIÓN JURÍDICA Y EFECTOS DEL INCUMPLIMIENTO
%
% Objetivos de esta sección:
%   1. Caracterizar el consentimiento informado (CI) como derecho fundamental del paciente
%      y deber jurídico del médico (art. 15 Ley 26842; arts. 4, 8 Ley 29414).
%   2. Analizar los requisitos del CI válido: información suficiente, comprensible,
%      oportuna; voluntad libre e informada; forma (escrita vs. verbal).
%   3. Distinguir el CI como:
%      - Eximente de responsabilidad (cuando es válido y previo al acto médico).
%      - Fuente autónoma de responsabilidad (cuando hay déficit informativo que
%        causa daño, aunque el acto médico sea técnicamente correcto).
%   4. Analizar el nexo causal en el CI: ¿habría el paciente rechazado el tratamiento
%      de haber sido informado adecuadamente? (test de causalidad hipotética).
%   5. Problemas especiales: urgencia médica, incapacidad del paciente, CI por representación.
%
% Referencias clave: Ley 29414 (arts. 4, 8, 15), Fernández Cruz, Galán Cortés
% Extensión estimada: 1.500–2.000 palabras

\section{El consentimiento informado: configuración jurídica y efectos del incumplimiento}
\label{sec:consentimiento}



% SECCIÓN 5: EL DAÑO MÉDICO — TIPOLOGÍA Y CRITERIOS JURISPRUDENCIALES DE VALORACIÓN
%
% Objetivos de esta sección:
%   1. Presentar la tipología del daño en el ámbito médico:
%      - Daño patrimonial: daño emergente (gastos médicos, rehabilitación) y
%        lucro cesante (pérdida de ingresos, incapacidad laboral).
%      - Daño extrapatrimonial: daño moral (sufrimiento psíquico) y daño a la persona
%        (daño al proyecto de vida, daño biológico) — distinción Fernández Sessarego.
%   2. Analizar el daño iatrogénico (causado por el tratamiento) y su diferencia
%      con el riesgo propio de la enfermedad (daño inevitable vs. evitable).
%   3. Examinar el nexo causal como elemento crítico:
%      - Teoría de la causa adecuada (art. 1985 CC).
%      - Causalidad concurrente (culpa del médico + riesgo propio de la enfermedad).
%      - Pérdida de chance: ¿es aplicable en el sistema peruano?
%   4. Criterios de cuantificación del daño en la jurisprudencia peruana:
%      ausencia de baremo, discrecionalidad judicial, estándares emergentes.
%
% Referencias clave: Espinoza Espinoza, Fernández Sessarego, art. 1985 CC, CAS. relevantes
% Extensión estimada: 1.500–2.000 palabras

\section{El daño médico: tipología y criterios jurisprudenciales de valoración}
\label{sec:dano_medico}



% SECCIÓN 6: ANÁLISIS DE LA JURISPRUDENCIA PERUANA — PATRONES Y CRITERIOS DOMINANTES
%
% Objetivos de esta sección:
%   1. Presentar el corpus de sentencias analizadas (Corte Suprema + Cortes Superiores,
%      período 2010–2024) y los criterios de selección.
%   2. Identificar los patrones jurisprudenciales en relación con:
%      a) Lex artis: ¿cómo la determinan los jueces? ¿recurren a peritos? ¿a protocolos?
%      b) Consentimiento informado: ¿cuándo se considera insuficiente? ¿qué consecuencias?
%      c) Daño médico: ¿qué tipos de daño se reconocen? ¿cómo se cuantifican?
%   3. Analizar la coherencia/incoherencia entre los criterios de la Sala Civil
%      Permanente y la Sala Civil Transitoria de la Corte Suprema.
%   4. Identificar los criterios determinantes del resultado (fundada/infundada):
%      ¿qué elemento es más frecuentemente deficiente — lex artis, nexo causal o daño?
%   5. Señalar las lagunas jurisprudenciales y propuestas de lege ferenda.
%
% ESTRUCTURA RECOMENDADA:
%   \subsection{Corpus y metodología}
%   \subsection{Criterios sobre la lex artis}
%   \subsection{Criterios sobre el consentimiento informado}
%   \subsection{Criterios sobre el daño médico y su valoración}
%   \subsection{Patrones sistémicos y vacíos jurisprudenciales}
%
% Extensión estimada: 2.500–3.500 palabras (sección central del artículo)

\section{Análisis de la jurisprudencia peruana: patrones y criterios dominantes}
\label{sec:jurisprudencia}



% SECCIÓN 7: CONCLUSIONES
%
% Objetivos de esta sección:
%   1. Sintetizar los hallazgos dogmáticos sobre la lex artis, el consentimiento informado
%      y el daño médico como ejes de la responsabilidad civil médica en el Perú.
%   2. Resumir los patrones jurisprudenciales identificados y sus implicancias prácticas.
%   3. Señalar las inconsistencias del sistema actual y las lagunas normativas.
%   4. Formular propuestas concretas de lege ferenda o de criterios jurisprudenciales
%      que refuercen la coherencia y predictibilidad del sistema.
%   5. (Opcional) Apuntar líneas de investigación futura.
%
% Formato: conclusiones numeradas (no subsecciones) — estilo habitual en revistas peruanas.
% Extensión estimada: 500–800 palabras

\section{Conclusiones}
\label{sec:conclusiones}



% ==========================================================
% REFERENCIAS
% ==========================================================

\clearpage
\small
\singlespacing
\printbibliography[title={Referencias}]

\end{document}
